% Part:sets-functions-relations
% Chapter: size-of-sets
% Section: comparing-sizes

\documentclass[../../../include/open-logic-section]{subfiles}

\begin{document}

\olfileid[fr]{sfr}{siz}{car}

\olsection{Ensembles de tailles différentes et théorème de Cantor}

\begin{explain}
Nous avons précisé ce que signifie, pour deux ensembles, avoir la
même taille. Nous pouvons aussi préciser ce que signifie être plus
petit qu'un autre ensemble. Pour définir « être plus petit ou
équipotent », nous demanderons une !!{injection} du premier ensemble
dans le second, au lieu d'une !!{bijection} entre les deux. S'il existe
une telle fonction, la taille du premier ensemble est inférieure ou
égale à celle du second. Intuitivement, une !!{injection} garantit
que l'image contient au moins autant d'éléments que le domaine,
puisque deux éléments distincts du domaine ne peuvent pas avoir
la même image.
\end{explain}

\begin{defn}
On dit que $A$ est \emph{de taille inférieure ou égale à celle
de}~$B$, et on écrit $\cardle{A}{B}$, si et seulement s'il existe
une !!{injection} $f \colon A \to B$.
\end{defn}

Cette relation est réflexive et transitive, mais elle n'est pas
symétrique ; on laisse la vérification en exercice. On peut aussi
définir une relation exprimant qu'un ensemble est strictement plus
petit qu'un autre.

\begin{defn}
On dit que $A$ est \emph{strictement plus petit que}~$B$, et on écrit
$\cardless{A}{B}$, si et seulement s'il existe une !!{injection}
$f\colon A \to B$ mais aucune !!{bijection} $g\colon A \to B$,
c'est-à-dire si $\cardle{A}{B}$ et $\cardneq{A}{B}$.
\end{defn}

Cette relation est irréflexive et transitive ; on laisse également
la vérification en exercice. Avec cette notation, $A$ est
!!{enumerable} si et seulement si $\cardle{A}{\Nat}$, et $A$ est
!!{nonenumerable} si et seulement si $\cardless{\Nat}{A}$.\footnote{
Précision éditoriale : la seconde caractérisation utilise une
hypothèse de choix laissée implicite dans l'original. Sous l'axiome
du choix, tout ensemble infini contient une partie équipotente à
$\Nat$, ce qui justifie l'injection de $\Nat$ dans~$A$. Sans
hypothèse de choix, la définition générale de la non-dénombrabilité
reste la négation de $\cardle{A}{\Nat}$ ; on ne peut pas en déduire
en général $\cardless{\Nat}{A}$.}
On peut ainsi reformuler
\oliflabeldef{sfr:siz:nen-alt:thm:nonenum-pownat}{%
le résultat du \olref[sfr][siz][nen-alt]{thm:nonenum-pownat}
par $\cardless{\Nat}{\Pow{\Nat}}$}{%
le résultat du \olref[sfr][siz][nen]{thm:nonenum-pownat}
par $\cardless{\PosInt}{\Pow{\PosInt}}$}.
En fait, \citet{Cantor1892} a montré que ce dernier phénomène est
\emph{tout à fait général} :

\begin{thm}[Cantor]\ollabel{thm:cantor}
Pour tout ensemble $A$, on a $\cardless{A}{\Pow{A}}$.
\end{thm}

\begin{proof}
L'application $f(x) = \{x\}$ est une !!{injection}
$f \colon A \to \Pow{A}$ : si $x \neq y$, alors
$\{x\} \neq \{y\}$ par extensionnalité, donc $f(x) \neq f(y)$.
Ainsi, $\cardle{A}{\Pow{A}}$.

\begin{editorial}
Si la \olref[nen]{sec} est présente, nous donnons la démonstration
détaillée ; sinon, nous donnons une démonstration plus brève,
correspondant à la \olref[nen-alt]{sec}.
Le lecteur complet conserve les deux démonstrations, dans cet ordre.
\end{editorial}

\frproofversions{Démonstration détaillée.}{Démonstration condensée.}{sfr:siz:nen:sec}{%
  Montrons maintenant qu'il ne peut exister de fonction
  !!{surjective} $g\colon A \to \Pow{A}$, et donc a fortiori
  aucune fonction !!{bijective}. Nous aurons ainsi
  $\cardneq{A}{\Pow{A}}$. Soit $g\colon A \to \Pow{A}$.
  Puisque $g$ est totale, chaque $x \in A$ a pour image une partie
  $g(x) \subseteq A$. Pour montrer que $g$ n'est pas surjective,
  définissons une partie $\overline{A} \subseteq A$ qui, par sa
  construction, ne peut pas appartenir à l'image de~$g$. Posons
  \[
  \overline{A} = \Setabs{x \in A}{x \notin g(x)}.
  \]
  Comme $g(x)$ est défini pour tout $x \in A$, cette définition
  donne bien une partie de~$A$. Pourtant, $\overline{A}$ ne peut
  pas appartenir à l'image de~$g$. Soit $x \in A$ quelconque ;
  montrons que $\overline{A} \neq g(x)$. Si $x \in g(x)$,
  il ne satisfait pas la condition $x \notin g(x)$, donc
  $x \notin \overline{A}$. Inversement, si
  $x \notin g(x)$, alors $x \in A$ et $x \notin g(x)$,
  donc $x \in \overline{A}$ ; tout $x \in \overline{A}$
  satisfait lui-même cette condition. Ainsi, pour chaque
  $x \in A$, on a $x \in g(x)$ si et seulement si
  $x \notin \overline{A}$, et donc $g(x) \neq \overline{A}$.\footnote{
  Note de traduction : l'original écrit ici « pour chaque
  $x \in \overline{A}$ ». Il faut considérer chaque $x \in A$
  pour exclure $\overline{A}$ de toute l'image de~$g$ ; les deux
  cas de l'argument sont explicités en conséquence.}
  Autrement dit, $\overline{A}$ n'appartient pas à l'image de~$g$,
  ce qui contredirait la surjectivité de~$g$.}{
  Il reste à montrer que $\cardneq{A}{\Pow{A}}$. Raisonnons par
  l'absurde : supposons que $\cardeq{A}{\Pow{A}}$, donc qu'il
  existe une !!{bijection} $g \colon A \to \Pow{A}$. Considérons
  \[
    D = \Setabs{x \in A}{x \notin g(x)}
  \]
  On a $D \subseteq A$, donc $D \in \Pow{A}$. Puisque $g$ est
  une !!{bijection}, il existe un $y \in A$ tel que $g(y) = D$.
  Mais alors :
  \[
    y \in g(y) \text{ si et seulement si } y \in D
    \text{ si et seulement si } y \notin g(y).
  \]
  C'est une contradiction ; on a donc $\cardneq{A}{\Pow{A}}$.}{}
\end{proof}

\begin{explain}
\oliflabeldef{sfr:siz:nen:thm:nonenum-pownat}{Comparons la démonstration
  du \olref{thm:cantor} à celle du
  \olref[nen]{thm:nonenum-pownat}. Nous y avons montré que, pour
  toute liste $Z_1$, $Z_2$, \dots{} de parties de~$\PosInt$,
  on peut construire un ensemble d'entiers~$\overline{Z}$ qui
  ne figure pas dans la liste. En effet, pour chaque
  $n \in \PosInt$, on a $n \in Z_n$ si et seulement si
  $n \notin \overline{Z}$. Il existe donc toujours un entier
  qui appartient à l'un des ensembles $Z_n$ et $\overline{Z}$,
  mais pas à l'autre. Nous reprenons ici la même idée : les
  indices~$n$ sont désormais des éléments de~$A$, au lieu
  d'être des éléments de~$\PosInt$. L'ensemble $\overline{A}$
  est défini de manière à différer de~$g(x)$ pour chaque
  $x \in A$, car $x \in g(x)$ si et seulement si
  $x \notin \overline{A}$. Là encore, pour chaque $x \in A$,
  cet élément appartient à l'un des ensembles $g(x)$ et
  $\overline{A}$, mais pas à l'autre. De même que
  $\overline{Z}$ ne peut pas figurer dans la liste $Z_1$, $Z_2$,
  \dots, $\overline{A}$ ne peut pas appartenir à l'image de~$g$.}{}

\oliflabeldef{sfr:siz:nen-alt:thm:nonenum-pownat}{Comparons la
  démonstration du \olref{thm:cantor} à celle du
  \olref[nen-alt]{thm:nonenum-pownat}. Nous y avons montré que,
  pour toute liste $N_0$, $N_1$, $N_2$, \dots{} de parties
  de~$\Nat$, on peut construire un ensemble d'entiers~$D$ qui
  ne figure pas dans la liste. En effet, pour chaque $n \in \Nat$,
  on a $n \in N_n$ si et seulement si $n \notin D$. Nous
  reprenons ici la même idée : les indices~$n$ sont désormais
  des éléments de~$A$ au lieu d'être des éléments de~$\Nat$.
  L'ensemble $D$ est défini de manière à différer de~$g(x)$
  pour chaque $x \in A$, car $x \in g(x)$ si et seulement si
  $x \notin D$.}{}

On peut aussi rapprocher cette démonstration de celle du paradoxe
de Russell, \olref[sfr][set][rus]{thm:russells-paradox}. Le
théorème de Cantor a en effet inspiré à Russell son propre paradoxe.
\end{explain}

\begin{prob}
  Montrer que, pour tout ensemble~$A$, il ne peut exister
  d'!!{injection} $g\colon \Pow{A} \to A$. Indication : supposons
  $g\colon \Pow{A} \to A$ !!{injective}. Considérer
  $D = \Setabs{g(B)}{B \subseteq A \text{ et } g(B) \notin B}$,
  puis poser $x = g(D)$. Utiliser l'injectivité de~$g$ pour obtenir
  une contradiction.
\end{prob}

\end{document}
